\documentclass{tufte-handout}

%\geometry{showframe}% for debugging purposes -- displays the margins

\usepackage{amsmath}

% Set up the images/graphics package
\usepackage{graphicx}
\setkeys{Gin}{width=\linewidth,totalheight=\textheight,keepaspectratio}
\graphicspath{{graphics/}}

\title{Some excercises}
\author[The Tufte-LaTeX Developers]{The Tufte-\LaTeX\ Developers}
\date{\today}  % if the \date{} command is left out, the current date will be used

% The following package makes prettier tables.  We're all about the bling!
\usepackage{booktabs}

% The units package provides nice, non-stacked fractions and better spacing
% for units.

\usepackage{amsmath,amsthm,amsfonts,amssymb,amscd,latexsym}
\usepackage{physics}
\usepackage{siunitx} 		% Permite utilizar el sistema internacional de unidades
\usepackage{mathptmx}
\usepackage{mathtools}
\usepackage{cancel}
\usepackage{centernot} 		% Notacion matematica
\usepackage{empheq}			% Recuadros para ecuaciones

% The fancyvrb package lets us customize the formatting of verbatim
% environments.  We use a slightly smaller font.
\usepackage{fancyvrb}
\fvset{fontsize=\normalsize}

% Small sections of multiple columns
\usepackage{multicol}

% Provides paragraphs of dummy text
\usepackage{lipsum}

% These commands are used to pretty-print LaTeX commands
\newcommand{\doccmd}[1]{\texttt{\textbackslash#1}}% command name -- adds backslash automatically
\newcommand{\docopt}[1]{\ensuremath{\langle}\textrm{\textit{#1}}\ensuremath{\rangle}}% optional command argument
\newcommand{\docarg}[1]{\textrm{\textit{#1}}}% (required) command argument
\newenvironment{docspec}{\begin{quote}\noindent}{\end{quote}}% command specification environment
\newcommand{\docenv}[1]{\textsf{#1}}% environment name
\newcommand{\docpkg}[1]{\texttt{#1}}% package name
\newcommand{\doccls}[1]{\texttt{#1}}% document class name
\newcommand{\docclsopt}[1]{\texttt{#1}}% document class option name

% For newtcbtheorem
\usepackage[most]{tcolorbox}
\usepackage{cleveref}
\usepackage{fancybox}		% Recuadros emph eqn

\usepackage{empheq}			% Recuadros para ecuaciones


\newtcbtheorem[]{prob}{Problem}%
    {enhanced,
    colback = black!5, %white,
    colbacktitle = black!5,
    coltitle = black,
    boxrule = 0pt,
    frame hidden,
    borderline west = {0.5mm}{0.0mm}{black},
    fonttitle = \bfseries\sffamily,
    breakable,
    before skip = 3ex,
    after skip = 3ex
}{def}



\begin{document}

\section{Exponential and logarithmic equations}

\begin{prob}{Change bsae formula}{pl-example1}
    Use the change base equation to give an exact numeric value using natural logarithm of $\log_{4}\qty[3e^7+2]$.
    
    Recalling the base change formula,
    \begin{gather}
        \log_a[x] = \frac{\log_b[x]}{\log_b[a]}.
    \end{gather}
    Now we identify that $a=4$, $b=e$ and $x=3e^7+2$.
    Replacing does values into the base change formula we get,
    \begin{align*}
        \log_{4}\qty[3e^7+2] &= \frac{\ln\left[3e^7+2\right]}{\ln\left[4\right]}
    \end{align*}


\begin{empheq}[box=\shadowbox]{equation*}
\log_{4}\qty[3e^7+2] = \frac{\ln\left[3e^7+2\right]}{\ln\left[4\right]}
\end{empheq}
\end{prob}

\begin{prob}{Solving exponential equations}{pl-example1}
    Solve the following equation \[3^{2x}-6\cdot3^x=0.\]

    We start by re-writting the expression with the help of the following property $a^{mn} = \left(a^{m}\right)^{n}$,
    \begin{gather*}
        \left(3^{x}\right)^2-6\cdot3^x=0.
    \end{gather*}
    Now, let's apply the following change of variable $a = 3^x$.
    Hence, we can write the equation as,
    \begin{gather*}
        a^2-6\cdot a=0.
    \end{gather*}
    Now we have a quadratic equation.
    We can apply the quadratic formula $a=-b\pm\sqrt{b^2-4ac}/(2a)$ to solve the quadratic equation.
    However, since $c=0$, we can solve it very simple as follows,
    \begin{align*}
        a^2-6\cdot a&=0 \\
        a^2&=6\cdot a \\
        \frac{1}{a}\cdot a^2&=\frac{1}{a}\cdot6\cdot a \\
        a&=6.
    \end{align*}
    Now that we have a numeric value for $a$, we can undo the change of variable and solve for $x$,
    \begin{align*}
        a &= 6 \\
        3^x &= 6 \\
        \log\left[3^x\right] &= \log\left[6\right] \\
        x\log\left[3\right] &= \log\left[6\right] \\
        x &= \frac{\log\left[6\right]}{\log\left[3\right]}
    \end{align*}

\begin{empheq}[box=\shadowbox]{equation*}
    x = \frac{\log\left[6\right]}{\log\left[3\right]}
\end{empheq}
\end{prob}

\marginpar{Note:
    See how the answer is actually the change base formula applied to $\log_3\left[6\right]=\frac{\log\left[6\right]}{\log\left[3\right]}$.
    This makes sense, since $\log_{3}[3^x]=x$.
    Therefore, if you choose to applied $\log_3$ instead of $\log$ you will get the same answer.
    However, to give a numeric value, you need to be able to use the base $3$ in your calculator.
    If you dont have that option, then use the change base formula answer.
}

\begin{prob}{Solving logarithmic equation}{label}
    Solve the following equation by finding the values of $x$ which the equation os true.
    \[\log_6\qty[2x+2] + \log_2\qty[2x-2] = 4.\]

    Let's start by reducing the logarithmic terms with the properties of logarithms and then solv for $x$.
    \begin{align*}
        \log_6\qty[2x+2] + \log_2\qty[2x-2] &= 4 \\
        \log_6\qty[\qty(2x+2)\qty(2x-2)] &= 4 \\
        6^{\log_6\qty[\qty(2x+2)\qty(2x-2)]} &= 6^{4} \\
        \qty(2x+2)\qty(2x-2) &= 6^{4} \\
        4x^2 -4x + 4x -4 &= 6^{4} \\
        4x^2 -4 &= 6^{4} \\
        4x^2 &= 6^{4}+4  \\
        x^2 &= \frac{6^{4}+4}{4}  \\
        x &= \pm 5\sqrt{13}.
    \end{align*}

    Now let's replace the numeric value into the original equation to see if both solutions are valid.
    We start with the positive one,
    \begin{align*}
        \log_6\qty[2\cdot5\sqrt{13}+2] + \log_2\qty[2\cdot5\sqrt{13}-2] &\approx \num{2.030990} + \num{1.969009} \\
                                                                        &\approx \num{3.999999}
    \end{align*}
    So this solution is valid.
    Now let's try the negative solution,
    \begin{align*}
        \log_6\qty[2\cdot-5\sqrt{13}+2] + \log_2\qty[2\cdot-5\sqrt{13}-2] &= \log_{6}\qty[\num{-34.05}] + \log_{6}\qty[\num{-38.05}] \\
                                                                          &= \mathrm{Math~error} + \mathrm{Math~error}.
    \end{align*}
    Since the logarithmic function is not defined for negativa numbers $x<0$, the second value is not a solution of the equation.

\begin{empheq}[box=\shadowbox]{equation*}
    x = 5\sqrt{13}
\end{empheq}


\end{prob}

\begin{prob}{Earthquakes}{labeln}
An earthquake of magnitude 8.5 was detected in the coastline of Michoacan. Recalling that the energy releases by an earthquake is $\log[E]=4.8+1.5 M$. Express the answers in tones of TNT (1tone of TNT = \num{4.184d9}).
\begin{itemize}
    \item What is the release energy in tones of TNT?
    \item What is the difference energy with an earthquake of magnitude 8.2?
\end{itemize}

    First we solve for the energy,
    \begin{align*}
        \log[E]=4.8+1.5 M\to E= 10^{4.8+1.5 M} = 10^{4.8}\cdot10^{1.5M}. 
    \end{align*}

    Now we can find the release energy in tones of TNT for an earthquake of magnitude \num{8.2}.
    \begin{align*}
        E_{8.5} &= 10^{4.8}\cdot10^{1.5\cdot8.5} \\
                &= 10^{4.8}\cdot10^{12.75} \\
                &\approx\num{3.548133d17}.
    \end{align*}
    Now let's express this number in tones of TNT by dividing the result over \num{4.184d9}.
    \begin{align*}
        E_{8.5} &\approx\frac{\approx\num{3.548133d17}}{\num{4.184d9}} \\
                &\approx\num{84802413.96}.
    \end{align*}

\begin{empheq}[box=\shadowbox]{equation*}
    E_{8.5}\approx\num{84802413.96}
\end{empheq}

    Now let's compute the difference energy between the earthquake of \num{8.5} with \num{8.2},
    \begin{align*}
        E_{8.5} - E_{8.2} &= 10^{4.8}\cdot10^{1.5\cdot8.5} - 10^{4.8}\cdot10^{1.5\cdot8.2} \\
                          &\approx \num{3.548133d17} - \num{1.258925d17} \\
                          &\approx\num{2.289208d17}.
    \end{align*}

    Now let's express the difference in terms of tones of TNT,
    \begin{align*}
        E_{8.5} - E_{8.2} &\approx\frac{\num{2.289208d17}}{\num{4.184d9}} \\
                          &\approx\num{54713384.32}
    \end{align*}

\begin{empheq}[box=\shadowbox]{equation*}
    E_{8.5} - E_{8.2} \approx\num{54713384.32}
\end{empheq}

    
\end{prob}

\paragraph{Exercise 1}
Suppose that you invest \$500 at a annual interest rate of 1.5\%, compounded monthly. 
\begin{itemize}
    \item How long will it take for your money to be \$10 000?
    \item How long will it take for your money to be \$10 000 if the rate is 3\%?
\end{itemize}

\noindent\makebox[\linewidth]{\rule{\textwidth}{0.4pt}}

\paragraph{Exercise 2}
The sound level $\beta$ in decibels (db), with and intensity I is given by $\beta(I)=10\log\qty[I/I_o]$ Where $I_o=\num{1d-12}$.
\begin{itemize}
    \item Calculate the sound level for a pneumatic drill that operates to and intensity of $I_o=\num{1d-12}$
    \item Calculate the intensity of the pain threshold of 120 db
\end{itemize}

\noindent\makebox[\linewidth]{\rule{\textwidth}{0.4pt}}

\paragraph{Exercise 3}
Use the change base equation to give and exact numeric fraction using natural logarithm of $\log_2[2e^2]$

\noindent\makebox[\linewidth]{\rule{\textwidth}{0.4pt}}

\paragraph{Exercise 4}
Solve the following equation $6^{x-3}=2^{x+2}$.





\end{document}
